\section{Introduction}
\label{sec:intro}

Large language model (LLM) agents are increasingly deployed as autonomous
software engineers~\cite{jimenez2024swebench,yang2024sweagent}, long-context
assistants~\cite{wu2024longmemeval}, and conversational
companions~\cite{maharana2024locomo}. Their effectiveness is largely
determined by \emph{harness engineering}: the design of the surrounding
code that orchestrates memory, retrieval, prompt routing, and tool
use~\cite{lee2026metaharness}. Automating this design through
propose--evaluate loops, where a second LLM (the \emph{proposer})
iteratively improves the harness, has emerged as a promising
paradigm~\cite{hu2024adas,khattab2023dspy}.

Existing approaches span a spectrum of information exposure.
ACE~\cite{zhang2025ace} and OpenEvolve~\cite{sharma2025openevolve}
compress iteration history into summaries or curated playbooks, feeding
the proposer a lossy but compact view.
Meta-Harness~\cite{lee2026metaharness} takes the opposite stance,
deploying a coding agent with full filesystem access to all prior
candidates, scores, and traces, on the premise that a capable agent's
native browsing and reasoning already constitute an effective search
gradient. This intuition is sound: a coding agent's chain-of-thought
\emph{is} the optimization gradient, and its ability to autonomously
navigate large workspaces makes it a natural proposer. Yet even powerful
proposers benefit from active management of which files they attend to.
As harnesses grow in complexity, unrestricted file access can dilute the
proposer's \textbf{attention} and diffuse its \textbf{optimization
direction}, analogous to how CODEFILTER~\cite{li2025codefilter} shows
that filtering retrieved context to only positive-impact chunks improves
code completion over providing all retrieved files. Explicit
mechanism-level guidance further enables \textbf{directed exploration},
steering the proposer toward structurally distinct alternatives rather
than repeatedly revisiting familiar axes such as prompt wording.

This paper introduces \textbf{Optimization Harness}, a modular framework
that wraps any propose--evaluate loop with three capabilities:
(i)~Docker file-scope sandboxing that restricts the proposer to a
curated workspace, (ii)~mechanism-level optimization-focus prompts that
steer exploration along named axes, and (iii)~labeled best/worst
reference iterations that provide imitation targets and anti-patterns.
Two adaptive context-selection policies build on this framework: a
\emph{progressive} policy that escalates a budget state machine
(\texttt{low}$\,\to\,$\texttt{medium}$\,\to\,$\texttt{high}) only after
a Pareto-frontier stall, and a \emph{bandit} policy that scores each
workspace file with a per-file UCB scorer and concentrates high-utility
files in a curated ``hot workspace'' each iteration.

Across three benchmarks (LoCoMo, LongMemEval, SWE-bench Verified) and
two proposer families (Claude Code with Kimi K2.6, and Codex GPT-5.4),
the progressive policy outperforms the fixed-high default in five of six
cells. The bandit policy achieves the single strongest result, lifting
SWE-bench Verified from 0.44 to 0.64 (+20\,pp) on the full 500-task
split. Across ten adaptive runs, 35 of 48 score-improvement events occur
at \texttt{low} or \texttt{medium} budget, confirming that
\emph{managing what the optimizer sees} is a complementary and effective
lever to improving the optimizer itself. Our contributions are:
\begin{enumerate}
    \item We formalize \emph{proposer context management} in harness
    optimization and show that curating file access improves attention
    focus, optimization direction, and exploration diversity.
    \item We introduce Optimization Harness, a modular framework with
    Docker sandboxing, mechanism-focus prompts, and reference labeling
    that enables pluggable context-selection policies.
    \item We propose two adaptive policies (progressive and bandit) and
    evaluate them on three benchmarks with two proposer families,
    demonstrating consistent improvements over the fixed-high baseline.
    \item We provide an empirical analysis of where breakthroughs occur
    across budget tiers, showing that curating attention is more
    efficient per opportunity than maximizing context.
\end{enumerate}
